
\section{结论}
\label{sec:conclusion}

我们已经实现了\cref{sec:outline-floer-homology}中所宣布的那个纲要的目标，证明了\trArnoldConjecture。

我们必须指出，除了这个猜想之外，\trFloerHomology{}及其各个亲缘理论还有许多对\trSymplectic{}几何与\textbf{\trContact}（contact）几何的应用。事实上，对于想要了解\trSymplectic\trTopology{}现代面貌的读者而言，本书的第二部分正可以作为一个出发点。这里为想要继续深入下去的读者给出一些方向与参考文献。

首先是对\trSymplectic\trManifold{}的\trLagrange{}\trSubmanifold{}的研究\footnote{我们想借此机会更新这份参考文献清单，特别补上该领域若干领军学者新近的著作与论文 [35,~71,~33,~10,~11]，以及本书作者之一所取得的新结果 [23,~24]。}，\trFloer{}在 [29]~中已经着手了这一研究；另可参见 Oh~[59,~60,~61]、Seidel~与\trFukaya{}~[34,~69,~70] 的工作，以及与我们更近的\trDamian{}~[22]~与 Gadbled~[36]~的工作。

其次，是我们的构造在带\trContact{}型边界的\trSymplectic\trManifold{}情形下的推广。在这种情形，\trFloerHomology{}不再像本书此处那样与\trMorseHomology{}\trIsomorphic，而是把边界的\textbf{\trClosedCharacteristic}（closed characteristic）考虑了进来。这就引出了由\trViterbo{}~[77]~发展起来的\trSymplectic\trHomology{}（另可参见 万恰~[58]~的文章，以及 齐利巴克与弗劳恩费尔德~[20]~的文章）。

另一个分支则通向\trContact\trManifold{}的研究。把\trFloer{}构造的一个变体用于\trContact\trManifold{}的\textbf{\trSymplectization}（symplectization），便给出一个由\trManifold{}的\trReeb{}\trVectorField{}的诸\trPeriodic\trTrajectory{}所生成的“\trContact{}\trHomology”。这个理论归功于\trEliashberg{}、吉文塔尔与\trHofer{}~[27]。另可参见 布尔乔亚~[15]~的工作。

不过，这就足以再写满另一本书了。
